\documentclass[conference]{IEEEtran}
\IEEEoverridecommandlockouts
\usepackage{cite}
\usepackage{amsmath,amssymb,amsfonts,amsthm}
\usepackage{graphicx}
\usepackage{booktabs}
\usepackage{multirow}
\usepackage{xcolor}
\usepackage{url}
\newtheorem{proposition}{Proposition}

\begin{document}

\title{Face Grouping That Keeps the User's Corrections:\\ Order-Invariant Incremental Clustering for\\ Local-First Photo Libraries}

% AUTHORS: replace the bracketed fields with the real details before submission (do NOT submit placeholders).
\author{\IEEEauthorblockN{Anugrah Singh}
\IEEEauthorblockA{\textit{Dept. of Computer Science and Engineering} \\
\textit{[College name]}\\
[City], India \\
[email]}
\and
\IEEEauthorblockN{Yadriksha Uprety}
\IEEEauthorblockA{\textit{Dept. of Computer Science and Engineering} \\
\textit{[College name]}\\
[City], India \\
[email]}
}

\maketitle

\begin{abstract}
Photo managers that group a user's photos by person on the user's own device must re-cluster whenever new photos
arrive. Two things then go wrong. Greedy incremental clustering makes the people found depend on the order in which
photos were imported, and re-clustering silently undoes the mistakes the user has already fixed. We describe the
face clusterer of Cove, a local-first media system, which addresses both. It adapts the core-point idea of incremental
DBSCAN to face embeddings, separates a strict link threshold from a looser attach threshold, and accepts the user's
corrections (``same person'', ``different people'') as must-link and cannot-link constraints that are applied on
every later re-clustering. Its output depends only on the set of faces and corrections, not on their order. We
evaluate it on 30 simulated personal libraries whose people never appear in the 30 tuning libraries, under eight
image conditions, with every baseline tuned on the same grid and data. With ten simulated corrections per library,
B-Cubed F1 rises from 0.984 to 0.994; after new photos arrive, all corrections still hold, against 27\% under plain
re-clustering. Without corrections, its accuracy is within about one point of the best batch baselines on degraded
photos, while leader-based methods reach an order agreement as low as 0.52. Profile faces remain hard (CPLFW F1 0.84
against 0.91 for average linkage), and a pose-aware similarity we tested did not fix this.
\end{abstract}

\begin{IEEEkeywords}
face clustering, constrained clustering, incremental clustering, order invariance, user corrections, local-first software, personal photo libraries
\end{IEEEkeywords}

% ---------------------------------------------------------------------------------------------------------------------
\section{Introduction}
People most often search a personal photo collection by person (``all photos of grandma''). Cloud services cluster
face embeddings on their servers; local-first managers such as Immich~\cite{immich}, PhotoPrism~\cite{photoprism},
FaceCluster~\cite{filonenko2025facecluster} and the Cove system described here do it on the user's own hardware, so
neither photos nor embeddings leave the device.

On a personal device the library is never complete. Photos arrive in batches over months and the people view must be
updated each time, and no clusterer is perfect, so users merge and split people by hand. This creates two problems
that a one-off benchmark does not show:
\begin{enumerate}
\item \textbf{Order dependence.} The cheapest update is a greedy \emph{leader} pass~\cite{hartigan1975clustering}: a new
face joins the most similar cluster above a threshold or starts a new one. The result depends on import order, so two
imports of the same photos can show different people. The first version of Cove used exactly this rule.
\item \textbf{Lost corrections.} If the system re-clusters from scratch after an import, the user's earlier fixes are
not part of its input, and the same mistakes come back. The user must fix them again after every import.
\end{enumerate}

Constrained clustering with must-link and cannot-link constraints~\cite{wagstaff2001copkmeans,davidson2005agglomerative},
its incremental form~\cite{davidson2007incremental} and interactive face annotation~\cite{tian2007face,cui2007easyalbum}
are all well established, and we do not claim these parts as new. Our contribution is a clusterer that combines them so
that both problems are solved at once, and an evaluation designed to measure exactly that:
\begin{itemize}
\item an incremental face clusterer whose output, including the user's corrections, is provably independent of the
order and batching of faces (Sec.~\ref{sec:method});
\item a measurement of whether corrections survive re-clustering when new photos arrive: with the clusterer keeping
them, 100\% hold; with plain re-clustering, 26--27\% (Sec.~\ref{sec:corr});
\item an identity-disjoint evaluation against seven baselines tuned on equal terms, under eight image conditions, with
paired bootstrap confidence intervals, which reports where our method is worse (Sec.~\ref{sec:acc});
\item a negative result for pose-aware similarity (Sec.~\ref{sec:pose}) and a CPU comparison of four lightweight recognisers (Sec.~\ref{sec:recog}).
\end{itemize}
Code, tuning grids and per-run results are released with the system.

% ---------------------------------------------------------------------------------------------------------------------
\section{Related Work}
\textbf{Face clustering.} Large-scale face clustering uses approximate rank-order distances~\cite{otto2018clustering}
or graph networks that predict links on a $k$-nearest-neighbour graph~\cite{wang2019linkage,yang2019affinity}. These
need labelled training clusters and usually a GPU, and they re-cluster the whole set. We compare against training-free
methods a local application can run on a CPU: leader clustering, a leader-plus-merge variant, single and average
linkage, DBSCAN~\cite{ester1996dbscan} and HDBSCAN~\cite{campello2013hdbscan}.

\textbf{Incremental and constrained clustering.} Incremental DBSCAN~\cite{ester1998incremental} showed that inserting a
point changes clusters only locally and gives the same result as re-running DBSCAN; our clusterer builds on this
core-point idea. OGMC~\cite{montero2022ogmc} clusters face streams with an online mixture model. Must-link and
cannot-link constraints were introduced for $k$-means~\cite{wagstaff2001copkmeans} and studied for agglomerative
clustering~\cite{davidson2005agglomerative}; Davidson \emph{et al.}~\cite{davidson2007incremental} update a clustering
as a user adds constraints over time. Constraints can also be used to tune a clusterer's thresholds~\cite{daniels2006threshold,vancraenendonck2017selection}.

\textbf{Interactive face grouping.} Photo album systems have long let users label or correct face groups and propagate
the labels~\cite{tian2007face,cui2007easyalbum,song2006context}, and FaceCluster~\cite{filonenko2025facecluster} is a
recent local, interactive photo organizer. In the sources we searched (SciSpace, Consensus and web search, abstracts
only) we did not find an evaluation of whether such corrections survive later re-clustering, or of order dependence
in personal libraries.

\textbf{Lightweight face recognition.} ArcFace~\cite{deng2019arcface} training with compact backbones such as
MobileFaceNets~\cite{chen2018mobilefacenets}, GhostFaceNets~\cite{alansari2023ghostfacenets} and
EdgeFace~\cite{george2024edgeface} gives accurate 512-D embeddings at a few hundred MFLOPs.

% ---------------------------------------------------------------------------------------------------------------------
\section{Method}\label{sec:method}
\subsection{Pipeline}
Cove detects faces with SCRFD-500M~\cite{guo2022scrfd} (InsightFace \texttt{buffalo\_s}, $320\times320$ input on CPU),
aligns the largest face of each photo to $112\times112$ with the standard five-point similarity transform, and embeds
it with EdgeFace-S~\cite{george2024edgeface} (Sec.~\ref{sec:recog}; MobileFaceNet \texttt{w600k\_mbf} is evaluated as
well). Embeddings $x_i\in\mathbb{R}^{512}$ are L2-normalised, so the inner product $s(i,j)=x_i^\top x_j$ is cosine
similarity, and stored in a FAISS flat inner-product index~\cite{johnson2021faiss}. All inference runs through ONNX
Runtime on the user's machine.

\subsection{Order-invariant incremental clustering}
Fix a link threshold $\tau_\ell$, an attach threshold $\tau_a\le\tau_\ell$, a neighbour count $m\ge1$ and a merge
threshold $\tau_g$. Let $N_\ell(i)=\{j\ne i: s(i,j)\ge\tau_\ell\}$.
\begin{itemize}
\item Face $i$ is \emph{core} if $|N_\ell(i)|\ge m$. Two core faces are \emph{linked} if $s(i,j)\ge\tau_\ell$, and
clusters are the connected components of the link graph over core faces.
\item A non-core face joins the cluster of its most similar core face with $s\ge\tau_a$; otherwise it stays alone.
Non-core faces never link two clusters, while the looser $\tau_a$ still recovers degraded faces of a known person.
\item Cluster centroids are merged by average-linkage agglomerative clustering at similarity $\tau_g$, which re-joins a
person split by a large appearance change.
\end{itemize}
At the tuned setting $m=1$, every face with at least one neighbour above $\tau_\ell$ is core, so the core step is
single linkage at $\tau_\ell$; the method is then single linkage, plus attachment at $\tau_a$, plus a centroid merge.
We present it as an adaptation of known parts, and Sec.~\ref{sec:acc} measures each part.

\textbf{Incremental update.} When a batch arrives, one FAISS range search at $\tau_a$ returns the new faces' neighbours
among all faces. Adjacency lists are extended, faces whose $|N_\ell|$ reaches $m$ become core, and each newly core face
is unioned (union-find) with its core neighbours above $\tau_\ell$. Core status and links only switch from false to
true as faces are added, so no earlier decision is undone. Attachment and the merge are recomputed from the stored
lists when labels are read; no new similarity search is needed.

\subsection{User corrections}
The user can state that faces $i$ and $j$ are the same person (must-link) or different people (cannot-link). The
corrections are stored with the clusterer and applied every time labels are computed, including after new photos:
\begin{enumerate}
\item must-links are applied first, in a fixed order;
\item core links are then added in decreasing similarity (ties by face id), and a union is skipped if it would put a
cannot-linked pair into one cluster, which builds a constrained maximum-spanning forest;
\item non-core faces attach to their most similar allowed core face;
\item the centroid merge is average linkage in which two clusters containing a cannot-linked pair may never merge.
\end{enumerate}
A cannot-link wins over a conflicting must-link. A user's ``these are different people'' therefore splits a wrong
merge at its weakest links, and ``same person'' joins two clusters, including all photos added to them later.

\begin{proposition}
For fixed thresholds and corrections, the labels after inserting a set $F$ of faces are the same for every insertion
order and every split of $F$ into batches (up to renaming of clusters and exact similarity ties).
\end{proposition}
\begin{proof}[Proof sketch]
After all insertions the adjacency lists hold exactly the pairs of $F$ with $s\ge\tau_a$, since each pair is found when
its later face arrives. Core status is a function of these lists, and every link between core faces is added, either
when the second face becomes core or when the later face arrives. Without corrections the union-find components are
therefore the connected components of the final link graph. With corrections, the components are built from the final
lists in an order fixed by similarity, so they too depend only on $F$ and the corrections. Attachment uses the final
lists, and the merge is a deterministic function of the components. Exactly equal real-valued similarities, where the
face-id tie-break (and hence the order) could matter, do not arise in practice; neither do contradictory corrections,
which are resolved in a fixed order.
\end{proof}

% ---------------------------------------------------------------------------------------------------------------------
\section{Experimental Protocol}\label{sec:protocol}
All rules below were fixed before any test number was computed.

\textbf{Identity-disjoint libraries.} LFW~\cite{huang2007lfw} identities are split once (seed 2026) into a tuning half
and a test half, separately for ``family'' identities ($\ge$10 photos; 79 per half) and ``strangers'' ($<$10 photos;
1{,}844 per half). A library draws 15--40 family people (up to 200 photos each) and 50--150 strangers (one photo each)
from its own half only: 30 tuning and 30 test libraries, with no person shared. \emph{Look-alike} test libraries
contain a test person and the 7--14 test people whose mean embeddings are most similar (a stand-in for relatives).
CALFW~\cite{zheng2017calfw} (cross-age) and CPLFW~\cite{zheng2018cplfw} (cross-pose) are LFW identities; we cluster
each whole, restricted to test-half identities (8{,}212 and 7{,}718 faces).

\textbf{Image conditions.} Every photo is re-embedded after degradation: \emph{low-res} (about 16\,px face width),
\emph{blur} (Gaussian, $\sigma=3$), \emph{JPEG} (quality 8), \emph{dark} (gamma 3 plus noise), \emph{occluded} (black
band over the eyes) and \emph{mixed} (one random condition per photo, like a phone library), plus \emph{clean}.

\textbf{Equal tuning.} Every method gets a grid over the same similarity range (0.20--0.70, step 0.05) and is tuned on
the same 30 tuning libraries $\times$ 7 conditions with the same rule: highest mean adjusted Rand index
(ARI)~\cite{hubert1985ari}. Grid sizes follow the number of parameters (from 11 for one-threshold methods to 360 for
ours). The recogniser was chosen on tuning data by a pre-declared rule: EdgeFace-S, the faster model, unless its tuning
ARI for our method is more than 0.005 below MBF's (it is 0.9822 vs.\ 0.9826, so EdgeFace-S is used). Tuned EdgeFace-S
setting: $\tau_\ell=0.65$, $m=1$, $\tau_a=0.50$, $\tau_g=0.55$.

\textbf{Test.} Order-dependent methods run in 5 shuffled orders, ours in 3 (an empirical check of the proposition). We
report B-Cubed precision, recall and F1~\cite{amigo2009bcubed}, and the \emph{order agreement}: the lowest ARI between
the partitions from two orders of the same library. Differences are tested with a paired bootstrap over libraries
(10{,}000 resamples, 95\% intervals).

\textbf{Corrections (Sec.~\ref{sec:corr}).} Each test library is split once into ``now'' (80\% of photos) and
``later'' (20\%, arriving afterwards). A simulated user makes up to ten corrections on the ``now'' clustering, each
time fixing the largest remaining mistake: a cluster mixing two people receives a cannot-link between the central face
of the main person and the most similar face of the other; a person split over two clusters receives a must-link
between the two central faces. The ``later'' photos then arrive and everything is clustered again, either with the
corrections kept by the clusterer or with plain re-clustering (the behaviour of the current app). Nothing in this
experiment is tuned.

\textbf{Hardware.} Intel Core i5-8250U laptop (4 cores, 1.6\,GHz), 12\,GB RAM, Windows~11, CPU only.

% ---------------------------------------------------------------------------------------------------------------------
\section{Results}
\subsection{Order dependence}\label{sec:order}
On the test libraries, the two leader-based methods produce different people for different import orders: First-Leader
reaches an order agreement of 0.56 and its F1 differs by up to 5.9 points between orders; adding a centroid merge
(Two-stage) still reaches 0.76 (Table~\ref{tab:acc}, last row). Our method gave identical partitions in every order and
batch split, with and without corrections, as the proposition predicts. Single and average linkage are also
invariant, and DBSCAN nearly so (0.997, from border faces), but these batch methods must be re-run on the whole
library after each import; HDBSCAN reached 0.98.

\begin{table*}[t]
\caption{B-Cubed F1 on identity-disjoint test data, EdgeFace-S embeddings (mean over 30 libraries; whole test-half sets
for CALFW/CPLFW). Every method tuned on the same grid range and tuning libraries. Bold: best. Last row: lowest order
agreement (ARI between two orders) over all library conditions.}
\label{tab:acc}
\centering
\setlength{\tabcolsep}{5pt}
\begin{tabular}{lcccccccc}
\toprule
Condition & First-Leader & Two-stage & Single link. & DBSCAN & HDBSCAN & Avg.\ linkage & Ours, no attach & \textbf{Ours} \\
\midrule
Clean      & 0.992 & 0.985 & 0.994 & 0.994 & 0.996 & 0.992 & 0.987 & \textbf{0.998} \\
Blur       & 0.980 & 0.991 & 0.986 & 0.986 & 0.995 & 0.993 & 0.990 & \textbf{0.997} \\
Dark       & 0.989 & 0.993 & 0.993 & 0.993 & 0.993 & 0.995 & 0.988 & \textbf{0.997} \\
JPEG q8    & 0.961 & \textbf{0.993} & 0.979 & 0.979 & 0.991 & \textbf{0.993} & 0.985 & \textbf{0.993} \\
Low-res    & 0.917 & \textbf{0.986} & 0.963 & 0.958 & 0.975 & \textbf{0.986} & 0.978 & 0.980 \\
Mixed      & 0.902 & \textbf{0.981} & 0.956 & 0.955 & 0.948 & \textbf{0.981} & 0.978 & 0.978 \\
Occluded   & 0.852 & \textbf{0.952} & 0.800 & 0.805 & 0.925 & \textbf{0.952} & 0.907 & 0.943 \\
Look-alike & 0.993 & 0.984 & 0.993 & 0.993 & 0.997 & 0.994 & 0.992 & \textbf{0.998} \\
CALFW (age)  & 0.681 & 0.902 & 0.945 & 0.793 & 0.358 & 0.940 & 0.919 & \textbf{0.947} \\
CPLFW (pose) & 0.600 & 0.897 & 0.846 & 0.516 & 0.520 & \textbf{0.909} & 0.892 & 0.841 \\
\midrule
Order agreement (min) & 0.562 & 0.758 & 1.000 & 0.997 & 0.983 & 1.000 & 1.000 & 1.000 \\
\bottomrule
\end{tabular}
\end{table*}

\subsection{Accuracy without corrections}\label{sec:acc}
Table~\ref{tab:acc} gives the result with equal tuning. Under these terms our method is no longer clearly ahead on
degraded photos, and we report it as such. Against average linkage, the strongest baseline, the paired difference is
within $\pm0.009$ F1 on every library condition: ours is significantly better on clean, blur and dark photos (+0.002
to +0.006), not significantly different on JPEG, occluded ($-0.009$, 95\% CI $[-0.020,+0.001]$) and look-alike
libraries, and significantly worse on low-res ($-0.006$, $[-0.008,-0.004]$) and mixed photos ($-0.003$,
$[-0.004,-0.001]$). Two-stage gives almost the same numbers, but is order-dependent. Against single linkage and
DBSCAN, ours is better by 0.4 to 14 points, most on occluded faces; this is the effect of the attach threshold, whose
removal (``Ours, no attach'') costs up to 3.6 points. With MBF embeddings the
pattern is the same, except that HDBSCAN is best on occluded faces (0.963 vs.\ 0.928).

On whole sets our method trades recall for precision. On CALFW it has the best F1 (0.947). On CPLFW it is 6.8 points
below average linkage (0.841 vs.\ 0.909, MBF: 0.787 vs.\ 0.916): its precision is higher (0.985 vs.\ 0.912), but it
splits many profile photos from the same person's frontal photos (recall 0.733 vs.\ 0.906). Average linkage, which
compares every face with every other, collects these weak but numerous links; our single-linkage core does not. The
gap therefore lies in the clustering rule as well as in the embedding. For a user, a split person costs one ``same
person'' correction, while a merged pair of people is harder to notice; this is why we kept the precise setting.

\subsection{Corrections that survive re-clustering}\label{sec:corr}
Table~\ref{tab:corr} shows the effect of corrections. Ten corrections per library raise F1 on the ``now'' photos from
0.984 to 0.994 for both recognisers (paired gain +0.010, CI $[+0.008,+0.011]$ for EdgeFace-S); the largest gains are on
occluded photos (0.946$\to$0.977) and mixed photos (0.975$\to$0.989). When the ``later'' photos arrive and the library
is clustered again, the clusterer keeps every correction (100\% still hold) and F1 stays at 0.992. Plain re-clustering,
which forgets them, returns to 0.985, and only 27\% of the corrections still hold (MBF: 26\%), so about three of every
four fixed mistakes come back. In every library, inserting all faces and corrections in three random orders gave
identical partitions (order agreement 1.000).

\begin{table}[t]
\caption{User corrections, EdgeFace-S, 30 test libraries. ``Now'': F1 after $k$ corrections. ``Later'': after 20\%
new photos and re-clustering, with corrections kept by the clusterer or lost (plain re-clustering).}
\label{tab:corr}
\centering
\setlength{\tabcolsep}{2.6pt}
\begin{tabular}{lccccccc}
\toprule
 & \multicolumn{4}{c}{Now: F1 after $k$ fixes} & \multicolumn{2}{c}{Later: F1} & Fixes \\
Condition & $k$=0 & 1 & 3 & 10 & lost & kept & held \\
\midrule
All (8 cond.) & 0.984 & 0.987 & 0.989 & 0.994 & 0.985 & 0.992 & 27\%\,/\,100\% \\
Low-res       & 0.974 & 0.975 & 0.978 & 0.986 & 0.980 & 0.987 & 30\%\,/\,100\% \\
Mixed         & 0.975 & 0.976 & 0.980 & 0.989 & 0.978 & 0.988 & 18\%\,/\,100\% \\
Occluded      & 0.946 & 0.959 & 0.966 & 0.977 & 0.943 & 0.969 & 12\%\,/\,100\% \\
\midrule
All, MBF      & 0.984 & 0.987 & 0.990 & 0.994 & 0.983 & 0.992 & 26\%\,/\,100\% \\
\bottomrule
\end{tabular}
\end{table}

\subsection{A negative result: pose-aware similarity}\label{sec:pose}
To target the CPLFW weakness, we gave frontal--profile pairs a similarity bonus
$s' = s + b\,\min(1, \lvert|y_i|-|y_j|\rvert/c)$, where $y$ is a yaw estimate from the detector's five landmarks.
$b$ and $c$ were tuned on the tuning libraries and on CFP-FP~\cite{sengupta2016cfp} frontal--profile pairs (a dataset
not used for testing), with a pre-declared rule to keep the bonus only if it improved the tuning score by at least
0.002. It passed this rule ($b=0.15$, the edge of the grid), but on the test data it failed: on CPLFW, precision fell
from 0.985 to 0.786 and F1 from 0.841 to 0.800, because the bonus joins profile photos of \emph{different} people as
readily as of the same person. With MBF, CPLFW F1 rose slightly (0.787$\to$0.796) only by losing precision, and
occluded F1 fell by 2.7 points. We did not adopt it.

\subsection{Choosing the face recogniser}\label{sec:recog}
We embedded identical crops with four ONNX models: \texttt{w600k\_mbf}, EdgeFace-S ($\gamma=0.5$) and EdgeFace-XS
($\gamma=0.6$)~\cite{george2024edgeface}, and GhostFaceNet W1.3~\cite{alansari2023ghostfacenets}. Table~\ref{tab:recog}
reports 10-fold verification accuracy on 3{,}000 same-person and 3{,}000 different-person pairs per set (CALFW,
CPLFW, AgeDB-30~\cite{moschoglou2017agedb}, CFP-FP~\cite{sengupta2016cfp}). For CALFW and
CPLFW these pairs are drawn at random from the images, \emph{not} the official hard-negative pair lists, so the numbers
are not comparable with published results. EdgeFace-S is never significantly worse than MBF and halves the embedding
time with one thread (2.1$\times$ faster; 1.67$\times$ with all cores). The final choice between the two was made by
the tuning rule of Sec.~\ref{sec:protocol}, not by these test sets.

\begin{table}[t]
\caption{Recognisers on identical crops: 10-fold verification accuracy (\%, random pairs; see text) and CPU latency
(ms/face). Brackets: difference to MBF; $^\ast$ = 95\% CI excludes zero.}
\label{tab:recog}
\centering
\setlength{\tabcolsep}{2.4pt}
\begin{tabular}{lcccc}
\toprule
 & MBF & EdgeFace-S & EdgeFace-XS & GhostFaceNet \\
\midrule
CALFW  & 96.40 & 96.57 (+0.17) & 96.28 ($-$0.12) & 96.62 (+0.22) \\
CPLFW  & 95.72 & 96.03 (+0.32) & 95.85 (+0.13) & 95.27 ($-$0.45$^\ast$) \\
AgeDB  & 96.15 & 96.92 (+0.77$^\ast$) & 95.97 ($-$0.18) & 97.40 (+1.25$^\ast$) \\
CFP-FP & 98.18 & 98.31 (+0.13) & 97.83 ($-$0.35) & 97.75 ($-$0.43$^\ast$) \\
\midrule
1 thread & 27.5 & \textbf{13.3} & \textbf{8.7} & 38.7 \\
All cores & 19.7 & \textbf{11.8} & \textbf{8.4} & 43.5 \\
Params (M) & 3.40 & 3.65 & 1.77 & 4.05 \\
\bottomrule
\end{tabular}
\end{table}

\subsection{Update time}\label{sec:time}
We added 8{,}290 clean LFW faces to a library in ten equal batches with the shipped EdgeFace-S setting, one thread.
After the last batch, the incremental update took 3.5\,s, against 4.7\,s for re-running our clusterer from scratch
and 16.7\,s for average linkage. The gain over re-running our own method is small (1.3$\times$): adding the batch
and updating links took only 0.15\,s, while the read-time steps (attachment and centroid merge) took 3.4\,s and grow
with the library. The update is therefore cheap but not constant-time. Successive partitions stayed consistent (ARI
$\ge0.998$ between the people found before and after each batch, on faces present in both).

% ---------------------------------------------------------------------------------------------------------------------
\section{Discussion and Limitations}
\textbf{What the results support.} A clusterer can give order-invariant people, accept the user's corrections and keep
them across later imports, at an accuracy close to the best batch baselines on degraded photos. The corrections result
is the main difference to re-clustering: without it, most fixed mistakes return after the next import.

\textbf{What they do not support.} (i) The libraries are simulations built from LFW celebrities; identities are
disjoint between tuning and test, but the photos come from one dataset, and real family albums (children ageing,
multi-face photos) were not tested. (ii) The simulated user always fixes the largest mistake and is never wrong; real
users fix fewer mistakes, in another order, and sometimes wrongly. (iii) Only the largest face per photo is clustered.
(iv) The tuned Two-stage threshold (0.70) lies on the edge of the grid. (v) Profile faces remain the weakest condition.
(vi) The corrections are implemented and tested in the clusterer; the current app does not yet persist them through
its re-clustering, which also overwrites user-assigned names. Wiring them in is the next engineering step.

\textbf{Future work.} Faces in one video track are the same person by construction, and a track often turns from
frontal to profile. In a simulation with ideal tracks on the CPLFW test half, adding a track's faces as must-links
raised same-person pair recall for the 171 people with a track from 0.585 to 0.784 at unchanged precision, but most of
this gain came from the extra faces themselves (0.719 unlinked), and overall F1 rose only from 0.854 to 0.860. Testing
this on real videos (e.g., YouTube Faces) is future work.

\textbf{Privacy.} All processing, including recognition, runs locally. We make no formal privacy claim beyond this
architecture. Stored face embeddings are as sensitive as the photos and should be protected by device encryption, and
the local API should accept requests only from the app itself (it currently accepts any origin).

\section{Conclusion}
Cove's face clusterer gives the same people for every import order and keeps the user's corrections across later
imports: with ten corrections per library, F1 rises from 0.984 to 0.994, and all corrections survive the next
re-clustering, against about a quarter with plain re-clustering. Without corrections it is within about one point of
the best batch baselines on degraded photos, and it remains weaker on profile faces, which a pose-aware similarity did
not fix.

\section*{Declarations}
\textbf{Data availability:} LFW, CALFW, CPLFW, AgeDB-30 and CFP-FP are public research datasets; code, splits, tuning
grids and per-run results are in the project repository (\texttt{evaluation/}, experiments 13--18).
\textbf{Ethics:} only public benchmark datasets were used; no new personal data were collected.
\textbf{Funding:} none. \textbf{Conflicts of interest:} none. \textbf{Author contributions:} to be completed by the
authors. \textbf{Use of AI tools:} an AI assistant (Claude, Anthropic) was used to write parts of the code and
experiment scripts, to search the literature and to edit the text. All numbers are produced by the released scripts
and were checked by the authors, who take responsibility for the content. \textbf{Licences:} InsightFace models and
EdgeFace weights are licensed for non-commercial research use.

\bibliographystyle{IEEEtran}
\bibliography{references}

\end{document}
